\documentclass[a4paper,fleqn]{cas-dc}

\usepackage[utf8]{inputenc}
\usepackage[T1]{fontenc}
\usepackage[expansion=false]{microtype}
\usepackage{amsmath,amssymb,amsfonts,amsthm}
\usepackage{algorithm}
\usepackage{algpseudocode}
\usepackage{graphicx}
\PassOptionsToPackage{compatibility=false}{caption}
\usepackage{subcaption}
\usepackage{booktabs}
\usepackage{multirow}
\usepackage{array}
\usepackage{float}
\usepackage{siunitx}
\usepackage{ragged2e}
\usepackage{placeins}
\usepackage[numbers,sort&compress]{natbib}
\usepackage{tikz}

\usetikzlibrary{
    arrows.meta,
    positioning,
    shapes.geometric,
    shapes.misc,
    calc
}

\newtheorem{corollary}{Corollary}[section]
\newtheorem{theorem}{Theorem}[section]
\newtheorem{lemma}{Lemma}[section]
\newtheorem{proposition}{Proposition}[section]
\newtheorem{remark}{Remark}[section]
\newtheorem{definition}{Definition}[section]
\newtheorem{example}{Example}[section]


\definecolor{revisionblue}{RGB}{0,0,180}
\definecolor{revisionred}{RGB}{180,0,0}
\newcommand{\RevisionAddBegin}{\begingroup\color{revisionblue}}
\newcommand{\RevisionAddEnd}{\endgroup}
\newenvironment{RevisionDeletion}{%
  \par\begingroup\color{revisionred}\footnotesize\ttfamily
  \setlength{\parindent}{0pt}\setlength{\parskip}{1pt}\obeyspaces
}{\par\endgroup}

\begin{document}
\let\WriteBookmarks\relax
\def\floatpagepagefraction{1}
\def\textpagefraction{.001}

\shorttitle{Net-Aware Learned Block Mapping for Binary-Image RDH}
\shortauthors{Ekodeck \textit{et al.}}

\title[mode=title]{Net-Aware Learned Block Mapping for Reversible Data Hiding in Binary Images}

\author[1,2,3]{St\'ephane Ga\"el R. Ekodeck}[
  orcid=0000-0002-8094-8832]
\cormark[1]
\ead{stephane-gael.ekodeck@facsciences-uy1.cm}
\credit{Conceptualization, Methodology, Software, Validation, Formal analysis, Investigation, Data curation, Visualization, Writing -- original draft}

\author[4,5]{Vivien Beyala Kamgang}[
  orcid=0000-0002-3644-1866]
\ead{vivsxx@univ-bertoua.cm}
\credit{Conceptualization, Methodology, Writing -- review and editing}

\author[1,2,3]{Serge Alain Eb\'el\'e}[
  orcid=0000-0001-7036-7068]
\ead{serge-alain.ebele@facsciences-uy1.cm}
\credit{Conceptualization, Methodology, Validation, Writing -- review and editing}

\author[5,4]{Julius Marcellin Nkenlifack}[
  orcid=0000-0003-3322-4687]
\ead{marcellin.nkenlifack@gmail.com}
\credit{Supervision, Project administration, Writing -- review and editing}

\affiliation[1]{
  organization={Department of Computer Science, LIRIMA, TEAM GRIMCAPE, University of Yaound\'e I},
  addressline={P.O. Box 812},
  city={Yaound\'e},
  country={Cameroon}
}
\affiliation[2]{
  organization={IRD, UMI 209 UMMISCO, IRD France Nord},
  city={Bondy},
  postcode={F-93143},
  country={France}
}
\affiliation[3]{
  organization={Sorbonne Universit\'e, UMI 209 UMMISCO},
  city={Paris},
  postcode={F-75005},
  country={France}
}
\affiliation[4]{
  organization={RTU Mathematics and Computer Science, The University of Bertoua},
  addressline={P.O. Box 416},
  city={Bertoua},
  country={Cameroon}
}
\affiliation[5]{
  organization={RUFCEA (UR-IFIA), University of Dschang},
  addressline={P.O. Box 96},
  city={Dschang},
  country={Cameroon}
}

\cortext[1]{Corresponding author}

\begin{abstract}
This paper presents a reversible data hiding scheme for plaintext binary
images under an explicit synchronization model. Its technical contribution is
serialization-aware net-capacity-driven activation of reversible PEAK/ZERO
$3\times3$ mappings: the active mapping configuration is selected by usable
mapping capacity after charging the exact serialized description required for
deterministic recovery. The implementation represents the active
PEAK subset by a combinatorial rank and the distance-one flip positions by
base-9 digits; a convolutional neural network (CNN) supplies the learned
distortion ranking for admissible ZERO candidates. On eight $512\times512$
\RevisionAddBegin
document images, this activation raises mean net payload from 2139 to 2245 bits
relative to an otherwise identical wire-charged gross-selection variant and
preserves exact round-trip reversibility. Frozen external validation on all 199
scanned FUNSD documents, with no training or tuning on FUNSD, preserves exact
message and cover recovery for every available run under both fixed-threshold
and Otsu binarization. Relative to all-feasible wire-charged activation, the
serialization-aware variant saves about 239--241 auxiliary bits per paired
image and is never worse in delivered net payload, although positive net
payload remains uncommon at targets up to 512 bits. In a common executable
protocol on 100 disjoint thresholded BOSSbase images, the proposed adaptive
block-mapping method (ABM), Dong, Huynh--Nguyen, and the opposite-center-pair
method (PPOCP) are compared with identical messages, payload targets,
distance-reciprocal distortion (DRD) computation, auxiliary-bit accounting,
and reversibility checks. The proposed method achieves the highest useful
capacity among the evaluated implementations, while PPOCP retains the
lowest-distortion and most favorable detectability profile. Logistic detector
AUC reaches 0.900 and 0.957 at 256 and 512 bits, placing the validated operating
region in capacity-oriented reversible annotation over a lossless channel with
detectability-aware payload selection.
\RevisionAddEnd
\end{abstract}

\begin{keywords}
\textcolor{revisionblue}{Binary-image reversible data hiding \sep PEAK/ZERO block mapping \sep}
\textcolor{revisionblue}{Serialization-aware mapping selection \sep Auxiliary-information coding \sep}
\textcolor{revisionblue}{Learned distortion ranking}
\end{keywords}

\maketitle

\section{Introduction}
\label{sec:introduction}

The proliferation of digital archives has increased the need to associate
authentication data, annotations, and integrity evidence with images while
preserving the original content. Cryptography protects message content
\RevisionAddBegin
\cite{Rivest1978RSA}, whereas steganography conceals its presence
\cite{Johnson2001InformationHiding}. Reversible data hiding (RDH) combines these objectives with
\RevisionAddEnd
exact cover restoration, a property that is especially valuable for medical,
legal, and administrative binary documents.

\RevisionAddBegin
Binary images have two pixel states, so every modification can alter
\RevisionAddEnd
stroke continuity, topology, or local readability. Specialized RDH methods
therefore rely on pattern statistics, overlapping neighborhoods, or reversible
block substitutions rather than direct grayscale-domain operations
\cite{Chhajed2023BlockDiagonal,Huynh2025BlockMapping,Yang2023Large,
Ren2023Sliding,Dong2020Adaptive,Zhang2019Magnification,Yin2020PPOCP,
Yin2020SymmetricFlipping,Ren2019Encrypted,Feng2015Secure,Xuan2008RunLength}. Their practical performance depends
on three coupled quantities: visual distortion, the side information required
for exact synchronization, and the statistical footprint visible to a
steganalyst.

\RevisionAddBegin
Existing methods occupy complementary parts of the capacity--distortion space.
Magnification-based embedding provides low distortion by increasing image size
\cite{Zhang2019Magnification}; PPOCP
\RevisionAddEnd
uses overlapping opposite-center pairs for perceptually selective changes
\cite{Yin2020PPOCP}; Huynh--Nguyen exploit the combinatorial capacity of
$3\times3$ block mapping \cite{Huynh2025BlockMapping}; and Dong et al.\ adapt
overlapping patterns to image statistics \cite{Dong2020Adaptive}. These
contributions motivate a unified evaluation in
which useful payload is measured after serialized metadata and all methods use
\RevisionAddBegin
the same messages, cover images, and distortion metric.
\RevisionAddEnd

\RevisionAddBegin
This work builds on established net-capacity accounting. The distinguishing
algorithmic decision is that the exact serialized description cost of the
reversible mapping configuration participates in deciding which PEAK/ZERO
tables remain active. A convolutional neural network ranks
distance-one ZERO candidates by predicted perceptual cost, while the
reversible mechanism remains a deterministic PEAK/ZERO substitution with an
enumerative auxiliary stream. Gaussian-process threshold exploration and the
Random Forest uniform-block agent are reported as diagnostic ablations; the
deployed distance-one codec consists of the deterministic PEAK/ZERO mapping,
CNN ranking, serialization-aware activation, and enumerative auxiliary stream.
\RevisionAddEnd

The main scientific contributions of this paper are:
\begin{enumerate}
\RevisionAddBegin
    \item Serialization-aware net-capacity-driven activation of reversible
    PEAK/ZERO mappings, where active tables are selected by usable mapping capacity
    after the exact serialized description cost required for deterministic
    recovery.
    \item An executable enumerative wire representation for the active PEAK
    subset and the associated distance-one flip positions, yielding a decoder
    contract $(I',A_{\mathrm{wire}})\mapsto(I,M)$.
    \item A learned candidate-cost ranker evaluated against Hamming-index,
    direct DRD-cost, transition-preserving, and connectivity-aware rankers, plus
    a CNN architecture ablation.
    \item A common executable protocol comparing the proposed adaptive
    block-mapping method (ABM), the opposite-center-pair method (PPOCP),
    Huynh--Nguyen, and Dong with identical
    covers, messages, payload targets, reversibility checks,
    distance-reciprocal distortion (DRD) implementation, auxiliary-bit
    accounting, paired statistics, resources, and grouped steganalysis.
    \item Frozen external validation of the deployed CNN-ranked codec on all
    199 real scanned FUNSD documents under fixed-128 and Otsu binarization,
    without retraining or tuning, including exact round-trip, availability,
    auxiliary-cost, net-payload, DRD, and paired A1/A2 activation analyses.
\RevisionAddEnd
\end{enumerate}

Section~\ref{sec:relatedwork} positions the work in binary RDH and
steganalysis. Section~\ref{sec:proposedmethod} defines the reversible mapping,
learned ranking, and net-capacity codec. Section~\ref{sec:experimental}
\RevisionAddBegin
describes the common protocol and reports document-image, BOSSbase, and FUNSD
results. Section~\ref{sec:discussion} synthesizes the practical operating
region, validity conditions, and remaining limitations, and
Section~\ref{sec:conclusion} concludes the study.
\RevisionAddEnd


\section{Related Work}
\label{sec:relatedwork}

\RevisionAddBegin
Binary-image RDH is constrained by the lack of grayscale redundancy: each
pixel flip can break strokes, isolate components, or change local topology.
The relevant literature therefore centers on pattern substitution, distortion
control, and the side information required to restore the cover exactly.
\RevisionAddEnd

\RevisionAddBegin
\subsection{Binary Pattern and Block Substitution}
\RevisionAddEnd

\RevisionAddBegin
Early binary-image RDH uses logical operations and run-length or pattern
statistics \cite{Xuan2008RunLength}. Ho et al.\ \cite{Ho2009PatternSubstitution} established
high-capacity reversible pattern substitution by exchanging frequent and
eligible replacement patterns. Dong et al.\ \cite{Dong2020Adaptive} refine this
line through adaptive overlapping patterns, PF/PFR handling, and a compressed
location map. Yin et al.\ \cite{Yin2020PPOCP} select opposite-center pattern
pairs (PPOCP), yielding a low-distortion profile. Zhang et al.\
\cite{Zhang2019Magnification} obtain very low distortion through image
magnification, but the carrier dimensions change. Huynh and Nguyen
\cite{Huynh2025BlockMapping} are the closest direct block-mapping reference:
they use non-overlapping $3\times3$ blocks, frequent PEAK patterns,
zero-frequency replacement patterns, and Hamming-distance-limited mappings.
These works establish pattern substitution, $3\times3$ PEAK/ZERO block
mapping, and Hamming-distance control as the relevant technical foundation.
\RevisionAddEnd

\RevisionAddBegin
Other binary-image variants expand the design space. Chhajed and Garg
\cite{Chhajed2023BlockDiagonal} use block-diagonal partition patterns, Ren
et al.\ \cite{Ren2023Sliding} use sliding windows, and Yang
\cite{Yang2023Large} prioritizes mixed black-white regions. They are useful
contextual comparators, with executable wire models left open for a direct
net-payload benchmark under the protocol used here.
\RevisionAddEnd

\RevisionAddBegin
\subsection{Distortion, Coding, and Reconstruction Information}
\RevisionAddEnd

\RevisionAddBegin
Binary document quality is measured with DRD because it accounts for spatially
non-uniform perceptual sensitivity \cite{Lu2004DRD}. Feng, Lu, and Sun
\cite{Feng2015Secure} further show that run continuity and connectivity are
important when choosing flippable pixels. These results motivate the ranker
ablations in Section~\ref{subsec:component-ablations}: Hamming ordering alone
is compared with direct DRD-cost, transition-preserving, connectivity-aware,
and CNN-based rankings.
\RevisionAddEnd

\RevisionAddBegin
Coding-based RDH for binary covers also predates this work. Zhang et al.\
\cite{Zhang2012OptimalCodes} improve binary-cover RDH schemes using optimal
codes, and Zhang et al.\ \cite{Zhang2015OptimalTransition} study optimal
transition probabilities under general distortion metrics. Xiao et al.\
\cite{Xiao2022BinaryCovers} propose a general-distortion RDH framework for
binary covers using reconstruction information and matrix embedding; Li et
al.\ \cite{Li2023MultiEmbedding} extend this line with multi-embedding.
These studies establish distortion-aware binary-cover coding and
reconstruction information as important conceptual prior art. Their
contribution is complementary to the plaintext $3\times3$ PEAK/ZERO
block-mapping codec developed here, where the serialized mapping description
itself determines which reversible mappings remain active.
\RevisionAddEnd

